% ===================================================== BÖLÜM 2
\section{Ana Makale: Wang vd. (2024) Ne Yaptı?}
\label{sec:wang}

\textbf{Künye:} Wang, W., Zhang, Z., Chen, Z., Wang, Y., Yu, C., \& Ma, B. 2024, ApJS, 270, 14,
“Long-term Orbital Period Variation of Hot Jupiters from Transiting Time Analysis using TESS Survey Data”
(arXiv:2310.17225; doi:10.3847/1538-4365/ad0847).

\subsection{Beş cümlede makale}
\begin{enumerate}[leftmargin=*]
\item Yörünge dönemi 10 günden kısa, kütlesi 0,3\,M$_\mathrm{J}$'den büyük ve TESS'ten önce de gözlenmiş
      \textbf{326 sıcak Jüpiter} seçtiler.
\item TESS'in 2 dakikalık ışık eğrilerinden \textbf{7068 geçişin} orta zamanını ölçtüler; bunları ağırlıklı olarak
      Ivshina \& Winn (2022) veritabanındaki eski (literatür) zamanlarla birleştirdiler.
\item Her gezegene iki model uydurdular: \textbf{doğrusal} (dönem sabit) ve \textbf{kuadratik} (dönem sabit hızla
      değişiyor). Aykırı noktaları 3$\sigma$ ve 5$\sigma$ kırpmayla attılar, ikisini de raporladılar.
\item Bir gezegeni “aday” saymak için iki şart koydular: (a) $|\dot P|$ sıfırdan en az 3$\sigma$ uzak,
      (b) $\Delta$BIC $>$ 10. Bu şartları sağlayan \textbf{26 aday} buldular: 18'inde dönem kısalıyor, 8'inde uzuyor.
\item Bu 26 adaya yeni bir dayanıklılık testi uyguladılar: \textbf{LOOCV}. \textbf{12 aday testi geçemedi}
      (“marjinal”), \textbf{14 aday geçti}. Olası fiziksel sebepler olarak gelgit, apsis presesyonu, Rømer etkisi,
      Applegate etkisi ve kütle kaybını tartıştılar ama hiçbir sistem için “şudur” demediler.
\end{enumerate}

\begin{basit}[Bizim projenin bu makaleye göre yeri]
Wang vd. adayları \emph{buldu} ama sebeplerini \emph{sistem sistem sınamadı}. Bizim işimiz: LOOCV'yi geçen
güvenilir adaylardan birkaçını alıp (1) güncel TESS verisiyle sinyal hâlâ duruyor mu bakmak, (2) duruyorsa
hocanın listesindeki 6 etkiden hangisinin sorumlu olduğunu belirlemek. Makalenin kendisi de bunu öneriyor:
“Rømer etkisi RV verisinin uydurulmasıyla doğrulanabilir” (Wang vd. 2024, Bölüm 5.3).
\end{basit}

\subsection{26 adayın tam listesi (Wang vd. 2024, Tablo 3)}

Aşağıdaki tablo makaledeki Tablo~3'ün 5$\sigma$ kırpmalı sütunlarıdır. N: toplam geçiş zamanı sayısı
(literatür + TESS), makalenin metninden.

\begin{center}
\small
\begin{longtable}{l r r r r r l}
\toprule
Gezegen & $\dot P$ (ms/yıl) & BIC$_\text{doğ}$ & BIC$_\text{kuad}$ & $\Delta$BIC & N (lit+TESS) & LOOCV\\
\midrule
\endhead
\multicolumn{7}{l}{\emph{Dönemi kısalanlar (18)}}\\
CoRoT-2\,b   & $-21{,}65\pm2{,}96$  & 634,56 & 579,03 & 55,5 & 118+5 & \textbf{geçti}\\
HD 189733\,b & $-8{,}44\pm2{,}44$   & 58,98  & 47,33  & 11,7 & 14+20 & kaldı\\
HAT-P-37\,b  & $-34{,}29\pm6{,}67$  & 221,71 & 195,71 & 26,0 & 27+8  & \textbf{geçti}\\
KELT-16\,b   & $-27{,}82\pm4{,}76$  & 103,63 & 70,77  & 32,9 & 46+65 & \textbf{geçti}\\
TrES-1\,b    & $-14{,}82\pm2{,}37$  & 355,98 & 321,61 & 34,4 & 41+34 & \textbf{geçti}\\
TrES-3\,b    & $-8{,}77\pm0{,}47$   & 1834,16& 1510,6 & 323,6& 190+94& kaldı\\
TrES-5\,b    & $-9{,}71\pm2{,}11$   & 216,66 & 194,11 & 22,6 & 55+66 & \textbf{geçti}\\
WASP-4\,b    & $-6{,}43\pm0{,}56$   & 342,39 & 213,62 & 128,8& 75+47 & \textbf{geçti}\\
WASP-10\,b   & $-27{,}74\pm3{,}49$  & 186,51 & 123,57 & 62,9 & 41+8  & \textbf{geçti}\\
WASP-12\,b   & $-30{,}19\pm0{,}92$  & 1256,13& 222,20 & 1033,9&107+82& \textbf{geçti}\\
WASP-16\,b   & $-51{,}98\pm12{,}87$ & 99,24  & 84,17  & 15,1 & 11+6  & kaldı\\
WASP-19\,b   & $-1{,}64\pm0{,}29$   & 693,77 & 661,02 & 32,8 & 87+55 & kaldı\\
WASP-22\,b   & $-71{,}76\pm17{,}12$ & 38,79  & 21,64  & 17,2 & 11+5  & \textbf{geçti}\\
WASP-45\,b   & $-169{,}21\pm21{,}09$& 350,59 & 289,88 & 60,7 & 10+14 & \textbf{geçti}\\
WASP-47\,b   & $-48{,}45\pm14{,}14$ & 77,18  & 65,36  & 11,8 & 27+4  & kaldı\\
WASP-80\,b   & $-23{,}04\pm5{,}15$  & 4017,15& 3996,25& 20,9 & 28+5  & kaldı\\
XO-3\,b      & $-31{,}59\pm5{,}24$  & 350,35 & 313,85 & 36,5 & 35+12 & kaldı\\
XO-4\,b      & $-62{,}57\pm17{,}48$ & 90,84  & 77,61  & 13,2 & 21+16 & \textbf{geçti}\\
\midrule
\multicolumn{7}{l}{\emph{Dönemi uzayanlar (8)}}\\
HAT-P-7\,b   & $+18{,}28\pm4{,}05$  & 112,96 & 92,35  & 20,6 & 11+66 & \textbf{geçti}\\
HAT-P-43\,b  & $+94{,}40\pm22{,}55$ & 70,24  & 53,24  & 17,0 & 7+41  & kaldı\\
HAT-P-44\,b  & $+119{,}67\pm25{,}89$& 63,69  & 42,49  & 21,2 & 13+14 & kaldı\\
WASP-1\,b    & $+22{,}50\pm5{,}23$  & 88,21  & 69,68  & 18,5 & 36+9  & \textbf{geçti}\\
WASP-6\,b    & $+20{,}66\pm5{,}19$  & 32,84  & 17,30  & 15,5 & 12+13 & kaldı\\
WASP-11\,b   & $+23{,}32\pm6{,}69$  & 102,65 & 89,70  & 13,0 & 24+11 & kaldı\\
WASP-17\,b   & $+77{,}64\pm8{,}19$  & 550,63 & 457,85 & 92,8 & 19+12 & kaldı\\
WASP-46\,b   & $+51{,}68\pm2{,}86$  & 2852,85& 2512,55& 340,3& 54+32 & \textbf{geçti}\\
\bottomrule
\end{longtable}
\end{center}

\textbf{LOOCV'yi geçen 14 aday} (hocanın “pass $\rightarrow$ kullan” dediği liste):
CoRoT-2\,b, HAT-P-37\,b, KELT-16\,b, TrES-1\,b, TrES-5\,b, WASP-4\,b, WASP-10\,b, WASP-12\,b, WASP-22\,b,
WASP-45\,b, XO-4\,b (11 kısalan) ve HAT-P-7\,b, WASP-1\,b, WASP-46\,b (3 uzayan).

% ===================================================== BÖLÜM 3
\section{LOOCV Testi: Ne Yapar, Geçmek ve Kalmak Ne Demek?}
\label{sec:loocv}

\begin{hoca}[Bu bölüm proje kapsamında değil]
Hoca bunu sana kişisel olarak “öğren” dedi. Projede LOOCV'yi kendimiz uygulamıyoruz; yalnızca Wang vd.'nin LOOCV
sonucunu (14 aday) seçimin başlangıç listesi olarak kullanıyoruz. Bu bölüm, hoca “LOOCV nedir, geçmek ne demek?”
diye sorduğunda cevap verebilmen için var. TÜBİTAK metninde (Kısım I) yöntem olarak geçmez.
\end{hoca}

\subsection{Tek cümlelik tanım}
\textbf{LOOCV (Leave-One-Out Cross-Validation, “birini dışarıda bırak çapraz doğrulama”)}: O$-$C verisindeki
noktaları \emph{tek tek} çıkarıp her seferinde modeli yeniden uydurarak, “dönem değişimi” sonucunun
\emph{tek bir ölçüme} bağlı olup olmadığını sınayan testtir.

\subsection{Günlük hayattan benzetme}
Bir sınıfın not ortalamasının “yükseliyor” olduğunu iddia ediyorsun. Ama acaba bu yükselişi tek bir öğrencinin
yanlış girilmiş 100'ü mü yaratıyor? Bunu anlamak için her öğrenciyi sırayla listeden çıkarıp ortalamayı yeniden
hesaplarsın. Hangi öğrenciyi çıkarırsan çıkar “yükseliyor” sonucu hâlâ duruyorsa sonucun sağlamdır.
Bir öğrenciyi çıkarınca sonuç kayboluyorsa, sonuç o öğrenciye bağlıdır ve şüphelidir. LOOCV tam olarak budur;
“öğrenci” yerine “geçiş zamanı”, “yükseliyor” yerine “dönem değişiyor” koy.

\subsection{Algoritma (Wang vd. 2024'teki haliyle)}
\begin{enumerate}[leftmargin=*]
\item Tüm $n$ geçiş zamanıyla doğrusal ve kuadratik modeli uydur; $\dot P$, $\sigma_{\dot P}$, $\Delta$BIC hesapla.
\item $i = 1$'den $n$'ye kadar: $i$'inci noktayı çıkar, kalan $n-1$ noktayla iki modeli yeniden uydur,
      yeni $\dot P_{(-i)}$, $\sigma_{(-i)}$, $\Delta\mathrm{BIC}_{(-i)}$ değerlerini kaydet.
\item Her $i$ için iki şartı kontrol et: \textbf{(a)} $|\dot P_{(-i)}|/\sigma_{(-i)} \ge 3$ ve
      \textbf{(b)} $\Delta\mathrm{BIC}_{(-i)} > 10$.
\item \textbf{GEÇTİ}: $n$ denemenin \emph{hepsinde} iki şart da sağlanıyor.
      \textbf{KALDI}: En az bir noktanın çıkarılması şartlardan birini bozuyor.
\end{enumerate}
Wang vd. bunu her aday için şekil olarak da çizmiştir (makaledeki Şekil~7 ve Şekil~A12–A35): üst panelde
her çıkarılan noktaya karşı $\dot P$, alt panelde $\Delta$BIC. Kodu Bölüm~\ref{sec:kod}'daki \texttt{loocv()}
fonksiyonudur (10 satır).

\subsection{Geçmek ne demek?}
\begin{itemize}[leftmargin=*]
\item \textbf{Sonuç tek bir ölçüme yaslanmıyor.} Sinyal, veri setinin genelinden geliyor.
\item Hangi geçiş zamanını atarsan at, dönem değişimi hâlâ 3$\sigma$ ve $\Delta$BIC $>$ 10 düzeyinde.
\item Bu yüzden bu adaylar, sebep araştırmasına \emph{değecek kadar} sağlam kabul edilir. Hocanın 14'ten
      seçmeni istemesinin sebebi budur: zaman harcayacağın gezegenin sinyalinin tek bir hatalı ölçümün ürünü
      olmaması gerekir.
\end{itemize}

\subsection{Kalmak ne demek?}
\begin{itemize}[leftmargin=*]
\item Sinyali \textbf{tek bir nokta} taşıyor. Çoğunlukla bu nokta, hata çubuğu gerçekçi olmayacak kadar küçük
      verilmiş eski bir yer gözlemidir. Örnekler (Wang vd. 2024, Bölüm 4 ve 5.1):
  \begin{itemize}
  \item \textbf{WASP-17\,b}: Bento vd. (2014)'ün tek bir ULTRACAM zamanı çıkarılınca $\dot P = +77{,}6$'dan
        $+3{,}4 \pm 4{,}5$\msyil'e düşüyor: sabit dönemle uyumlu.
  \item \textbf{XO-3\,b}: Tek bir Spitzer zamanı çıkarılınca $-31{,}6$'dan $-8{,}8\pm5{,}1$'e ve
        $\Delta$BIC = 2,4'e iniyor.
  \item \textbf{TrES-3\,b}: Hatası 1 saniye (!) olarak verilmiş bir nokta çıkarılınca $\dot P$'nin işareti
        bile değişiyor.
  \end{itemize}
\item Kalmak “dönem değişmiyor” demek \emph{değildir}; “bu veriyle güvenle söyleyemeyiz” demektir
      (Wang vd. bunlara “marjinal aday” der).
\end{itemize}

\subsection{LOOCV'nin söylemediği şeyler (jüri sorarsa)}
\begin{dikkat}[Geçmek, “yörünge gerçekten bozuluyor” demek DEĞİLDİR]
\begin{enumerate}[leftmargin=*]
\item \textbf{LOOCV fiziksel sebebi söylemez.} WASP-4\,b testi geçer ama dönem değişiminin büyük kısmı büyük
      olasılıkla Rømer etkisidir (Bouma vd. 2020; Winn \& Stefánsson 2025). Sebep için Bölüm~\ref{sec:etkiler}–\ref{sec:karar}.
\item \textbf{LOOCV yalnızca \emph{tek} noktaya karşı korur.} Aynı makaleden gelen 9 noktanın hepsi aynı zaman
      sistemi hatasını taşıyorsa (HAT-P-7 örneği, Bölüm~\ref{sec:kavram}) LOOCV bunu yakalayamaz. Bunun çaresi
      kaynak kaynak kontrol ve “bu makaleyi çıkarınca ne oluyor?” sorusudur.
\item \textbf{LOOCV veri setine bağlıdır.} Yeni veri gelince sonuç değişebilir: TrES-5\,b Wang'da testi geçer, ama
      Rothmeier vd. (2025) yeni verilerle $\dot P$'nin $-5{,}3\pm2{,}2$\msyil'e düştüğünü ve tek bir erken nokta
      çıkarılınca tamamen kaybolduğunu gösterdi.
\end{enumerate}
\end{dikkat}

\subsection{Kendi yeniden üretimim: açık katalog verisiyle 14 aday}
Wang vd.'nin LOOCV'yi geçen 14 adayını, \textbf{yalnızca herkesin indirebileceği iki tabloyla} (Ivshina \& Winn
2022 literatür+TESS zamanları ve Wang vd. 2024 TESS zamanları; ETD eklemeden) Bölüm~\ref{sec:kod}'daki kodla
yeniden analiz ettim (5$\sigma$ kırpma, aynı iki şart, aynı LOOCV):

\begin{center}
\small
\begin{tabular}{l r r r r l}
\toprule
Gezegen & $n$ (TESS) & $\dot P$ (ms/yıl) & $|\dot P|/\sigma$ & $\Delta$BIC & LOOCV (benim)\\
\midrule
CoRoT-2\,b  & 122 (5)  & $-49{,}18\pm2{,}34$  & 21,0 & 435,8 & geçti\\
HAT-P-37\,b & 17 (8)   & $-17{,}87\pm16{,}95$ & 1,1  & $-1{,}7$ & \textbf{kaldı}\\
KELT-16\,b  & 89 (65)  & $-22{,}97\pm6{,}70$  & 3,4  & 7,2   & \textbf{kaldı}\\
TrES-1\,b   & 74 (34)  & $-18{,}99\pm2{,}50$  & 7,6  & 53,2  & geçti\\
TrES-5\,b   & 121 (66) & $-9{,}72\pm2{,}06$   & 4,7  & 17,5  & \textbf{kaldı}\\
WASP-4\,b   & 122 (47) & $-7{,}54\pm0{,}73$   & 10,3 & 101,0 & geçti\\
WASP-10\,b  & 49 (8)   & $-27{,}67\pm3{,}48$  & 8,0  & 59,2  & geçti\\
WASP-12\,b  & 189 (82) & $-28{,}64\pm0{,}98$  & 29,2 & 854,6 & geçti\\
WASP-22\,b  & 16 (5)   & $-75{,}08\pm18{,}84$ & 4,0  & 13,1  & \textbf{kaldı}\\
WASP-45\,b  & 24 (14)  & $-291{,}0\pm27{,}9$  & 10,4 & 105,3 & geçti\\
XO-4\,b     & 36 (16)  & $-58{,}74\pm16{,}97$ & 3,5  & 8,4   & \textbf{kaldı}\\
HAT-P-7\,b  & 77 (66)  & $+22{,}59\pm5{,}05$  & 4,5  & 15,6  & geçti\\
WASP-1\,b   & 45 (16)  & $+22{,}13\pm5{,}07$  & 4,4  & 15,3  & \textbf{kaldı}\\
WASP-46\,b  & 87 (33)  & $+65{,}20\pm3{,}53$  & 18,5 & 336,9 & geçti\\
\bottomrule
\end{tabular}
\end{center}

\textbf{Ne öğrendik?} Wang'ın “geçti” dediği 14 adaydan 6'sı (HAT-P-37, KELT-16, TrES-5, WASP-22, XO-4, WASP-1)
açık katalog verisiyle şartları sağlamıyor. Sebep çoğunlukla Wang'ın ETD'den eklediği ek noktalardır: bu
sistemlerde sinyal az sayıdaki amatör/eski veriye dayanıyor. CoRoT-2\,b'de ise benim değerim ($-49$) Wang'ınkinden
($-22$) çok farklı, çünkü Wang bazı literatür hata çubuklarını elle büyüttü; Adams vd. (2024) de CoRoT-2\,b'nin
literatür zamanlarında birden fazla hata buldu. Bu tablo seçimde ölçüt olarak kullanılmadı; ama seçilen dört hedefin \emph{açık veriyle de} sağlam
çıktığını gösteren \textbf{destekleyici kanıttır} (jüri “sonucu yeniden ürettin mi?” diye sorarsa).

% ===================================================== BÖLÜM 4
\section{14 Adaydan Hedef Seçimi: Gerekçe}
\label{sec:secim}

\subsection{Hocanın istediği ve bizim yaklaşımımız}
Hoca: “14 tanesine bakmak zor, birkaçını seç; ama kafana göre seçersen sorun olur, bir gerekçe bul. Parametre
tablosuna bak, belki ondan bir şey bulursun.” Yani \textbf{kaç tane} seçileceği de \textbf{neye göre} seçileceği de bize
bırakılmış; tek şart gerekçenin sağlam olması.

Jüri karşısında en sağlam gerekçe şudur: \textbf{sayıyı biz seçmeyiz; önceden yazılmış ölçütler seçer.} Ölçütleri
sonuçlara bakmadan yazarsın, 14 adaya uygularsın, kaç tane kalırsa o kadar hedefin olur. Kimse “neden 4? neden bu
4?” diyemez, çünkü cevap “ölçütleri sağlayan bunlar” olur. Ölçütlerin her biri de \textbf{hocanın bir cümlesine}
dayanır.

\subsection{Üç ölçüt}
\begin{center}
\begin{tabularx}{\textwidth}{c >{\raggedright\arraybackslash}X >{\raggedright\arraybackslash}X >{\raggedright\arraybackslash}X}
\toprule
 & Ölçüt & Nereden ölçülüyor? & Hocanın hangi sözüne dayanıyor?\\
\midrule
K1 & LOOCV'yi geçmiş & Wang vd. 2024, Tablo 3 & “LOOCV'yi geçen 14 tane”\\
K2 & TESS geçiş zamanı sayısı $\ge$ 30 & Wang vd. 2024, her adayın metninde verilen TESS sayısı (\textbf{makalenin parametresi}) &
  “TESS verilerini dahil edeceksin.” 30'un altında TESS katkısı tipik olarak tek bir sektördür ($\sim$27 gün); TESS ile kontrol yapılamaz.\\
K3 & Arşiv RV'sinden öngörülen Rømer katkısının belirsizliği $<$ gözlenen $|\dot P|$'nin üçte biri & DACE, Knutson vd. 2014,
  Bouma vd. 2020, Bonomo vd. 2017 (yayımlanmış ölçümler) & “Hangisinden etkilendiğini tespit et.” Gelgit ile Rømer yalnız
  O$-$C'den ayrılamaz (Bölüm~\ref{sec:rv}); RV yeterince hassas değilse o hedefte hocanın son sorusu cevapsız kalır.\\
\bottomrule
\end{tabularx}
\end{center}

\textbf{K3 neden “üçte bir”?} Gözlenen $\dot P$ ile Rømer öngörüsü arasında 3$\sigma$'lık bir farkı görebilmek için
öngörünün belirsizliği, ölçülen sinyalin üçte birinden küçük olmalıdır. Bu, projenin her yerinde kullandığımız
“3$\sigma$” kuralının aynısıdır; yeni ve keyfi bir eşik değildir.

\subsection{Huniyi uygulayalım}
\textbf{K1 $\rightarrow$ 14 aday.}

\textbf{K2 $\rightarrow$ 7 aday kalır.} TESS sayısı (Wang vd.): CoRoT-2 (5), HAT-P-37 (8), WASP-10 (8), WASP-22 (5),
WASP-45 (14), XO-4 (16), WASP-1 (9) elenir. Kalanlar: KELT-16 (65), TrES-1 (34), TrES-5 (66), WASP-4 (47), WASP-12 (82),
HAT-P-7 (66), WASP-46 (32).

\textbf{K3 $\rightarrow$ 4 aday kalır.} Her aday için bütün arşivlere baktım (DACE, HARPS RVBank, Knutson 2014,
Bouma 2020, Bonomo 2017) ve en iyi RV setiyle Rømer öngörüsünün belirsizliğini \texttt{rv\_analiz.py} ile hesapladım:

\begin{center}
\small
\begin{tabularx}{\textwidth}{l >{\raggedright\arraybackslash}X r r r l}
\toprule
Aday & En iyi RV verisi & $\sigma_\text{Rømer}$ & $|\dot P|$ (Wang) & $|\dot P|/3$ & K3\\
\midrule
WASP-12\,b & DACE: 90 RV (SOPHIE, ESPRESSO), 13,9 yıl & 1,7 & 30,2 & 10,1 & \textbf{sağlıyor}\\
WASP-4\,b  & Bouma 2020: 48 RV (CORALIE, HARPS, HIRES), 11,9 yıl & 0,3 & 6,4 & 2,1 & \textbf{sağlıyor}\\
HAT-P-7\,b & Knutson 2014: 43 RV (HIRES), 5,9 yıl & 1,2 & 18,3 & 6,1 & \textbf{sağlıyor}\\
TrES-1\,b  & Bonomo 2017: 15 RV (HARPS-N), 2,8 yıl & 2,2 & 14,8 & 4,9 & \textbf{sağlıyor}\\
KELT-16\,b & Hiçbir arşivde yok (DACE, HARPS, Bonomo) & -- & 27,8 & 9,3 & sağlamıyor\\
TrES-5\,b  & DACE: 8 RV (TRES), 0,6 yıl & 29 & 9,7 & 3,2 & sağlamıyor\\
WASP-46\,b & DACE: 21 RV (CORALIE, HARPS) & 56 & 51,7 & 17,2 & sağlamıyor\\
\bottomrule
\end{tabularx}
\end{center}
($\sigma_\text{Rømer}$ ve $\dot P$ ms/yıl. “RV var” ile “RV işe yarıyor” farklıdır: WASP-46'nın 21 RV'si var ama iki
farklı aletin sıfır noktaları ayrı uydurulunca eğim bilgisi kalmıyor.)

\textbf{Sonuç: WASP-12\,b, WASP-4\,b, HAT-P-7\,b, TrES-1\,b.}

\begin{basit}[Elenen ama RV'si iyi olanlar (yedek)]
K2'de elenen adaylardan WASP-10 ($\sigma_\text{Rømer} = 3{,}9$), WASP-22 (1,4) ve WASP-1 (2,6) iyi RV verisine sahip; ama TESS
verileri çok az (5–9 geçiş). Bir hedef proje sırasında sorun çıkarırsa yerine \textbf{WASP-10\,b} alınır (Kısım I, risk tablosu).
\end{basit}

\subsection{Bu dört hedef bir arada neden iyi?}
Ölçütler hedefleri seçti; ama sonuç ayrıca dört farklı fiziksel durumu temsil ediyor. Bu, hocanın 6 etkisinin hepsinin en
az bir hedefte gerçekten sınanmasını sağlar:
\begin{center}
\small
\begin{tabularx}{\textwidth}{l >{\raggedright\arraybackslash}X >{\raggedright\arraybackslash}X}
\toprule
Hedef & Rolü & Sınanacak ana soru\\
\midrule
\textbf{WASP-12\,b} & \textbf{Kalibrasyon.} Gelgit bozulması doğrulanmış tek sistem (Yee vd. 2020; Winn \& Stefánsson 2025). &
  Yöntemim bilinen sonucu ($\approx -29$\msyil) buluyor mu? Rømer ve apsis dışlanıyor mu?\\
\textbf{WASP-4\,b} & \textbf{Tartışmalı vaka.} Baştürk vd. (2025) “bozuluyor”; Bouma vd. (2020) ve Winn \& Stefánsson (2025)
  “ışık-yolu etkisi” diyor. & Arşiv RV'sinden öngörülen Rømer değeri $\dot P$'yi açıklıyor mu?\\
\textbf{HAT-P-7\,b} & \textbf{Dönemi uzayan sistem.} Gelgit (yavaş dönen yıldızda) dönemi uzatamaz. Bilinen eşlikçiler var. &
  Dönem artışı eşlikçinin ivmesinden mi geliyor?\\
\textbf{TrES-1\,b} & \textbf{Ne gelgit ne Rømer.} Gelgit için Q$'_\star \approx 8\times10^2$ gerekiyor (kuramdan $\sim$100 kat küçük);
  HARPS-N RV'leri Rømer'i dışlıyor. & Geriye kalan apsis, dinamik 3. cisim ve Applegate'ten hangisi?\\
\bottomrule
\end{tabularx}
\end{center}

\begin{hoca}[Hocaya / jüriye söylenecek hazır gerekçe]
“Wang vd. (2024)'nin LOOCV testini geçen 14 adaya sonuçlara bakmadan önce iki ölçüt daha uyguladık. Birincisi
makalenin kendi parametresi: en az 30 TESS geçiş zamanı. Projenin amacı TESS verisiyle kontrol olduğu için bu,
7 aday bıraktı. İkincisi, hangi fiziksel etkinin sorumlu olduğunu ayırt edebilmek için gereken arşiv radyal hız
verisinin yeterliliği: Rømer etkisinin öngörüsü, gözlenen dönem değişiminin üçte birinden daha hassas olmalı.
Bu iki ölçütü WASP-12\,b, WASP-4\,b, HAT-P-7\,b ve TrES-1\,b sağladı. Bu dört hedef sırasıyla doğrulanmış bir gelgit
bozulmasını, tartışmalı bir gelgit–ivme vakasını, dönemi uzayan bir sistemi ve gelgitle de ivmeyle de açıklanamayan
bir kısalmayı temsil ediyor.”
\end{hoca}

\subsection{Hoca “daha az” derse}
Hedef sayısı ölçütlerden çıktığı için 4'te kalmak en sağlamıdır. Hoca yine de 3 isterse, eleme için Wang tablosundaki
\textbf{sinyalin anlamlılığını} ($|\dot P|/\sigma$) kullan; en zayıf sinyal elenir:
WASP-12 (32,8), WASP-4 (11,5), TrES-1 (6,3), HAT-P-7 (4,5 $\rightarrow$ çıkar). Bedeli: dönemi \emph{uzayan} tek örnek gider.

\textbf{Neden $\Delta$BIC'e göre sıralamıyoruz?} $\Delta$BIC veri sayısıyla büyür (WASP-12'de 189, TrES-1'de 75 nokta) ve
iyimser hata çubuklarıyla şişer; Wang vd. de tek bir aykırı noktanın $\Delta$BIC'i ciddi biçimde değiştirebildiğini
yazar. Yani büyük $\Delta$BIC “daha güvenilir sinyal” demek değildir; jüri bunu kolayca sorgular. $|\dot P|/\sigma$ ise
projenin her yerinde kullandığımız “3$\sigma$” diliyle aynıdır.

\subsection{“Neden şunu seçmedin?” sorularına hazır cevaplar}
\begin{itemize}[leftmargin=*]
\item \textbf{KELT-16\,b?} TESS verisi bol ama hiçbir arşivde RV'si yok (K3). Ayrıca sonraki çalışmalar dönem değişimi bulamadı
      (Kutluay vd. 2026'nın literatür özeti); açık katalog verisiyle bende de $\Delta$BIC = 7,2.
\item \textbf{TrES-5\,b?} RV'leri yalnızca 0,6 yıla yayılıyor (K3); Rothmeier vd. (2025) sinyalin yeni veriyle kaybolduğunu gösterdi.
\item \textbf{WASP-46\,b?} RV'leri eğim ölçemeyecek kadar belirsiz (K3); Wang vd. bile artıklardaki saçılmanın kuadratik modele
      şüphe düşürdüğünü yazar.
\item \textbf{CoRoT-2, HAT-P-37, WASP-10, WASP-22, WASP-45, XO-4, WASP-1?} TESS geçiş sayısı 5–16 (K2).
\end{itemize}

\subsection{Destekleyici fiziksel tablo (14 aday)}
Aşağıdaki sayılar NASA Exoplanet Archive (\texttt{pscomppars}, Kasım 2026) parametrelerinden kendi hesabımdır.
“Gelgit gücü” $(M_p/M_\star)(R_\star/a)^5$'tir (Bölüm~\ref{sec:etkiler}); “Q$'_\star$ gereken”, Wang'ın $\dot P$'sini
gelgitle açıklamak için gereken değerdir (makul aralık $\sim10^5$–$10^7$). “RV eğimi gereken”, aynı $\dot P$'yi
Rømer ile açıklamak için yıldızın RV'sinin ne kadar hızla değişmesi gerektiğidir.

\begin{center}
\small
\begin{tabular}{l r r r r r r}
\toprule
Gezegen & $P$ (gün) & $V$ & $R_\star/a$ & Gelgit gücü ($\times10^{-8}$) & Q$'_\star$ gereken & RV eğimi gereken (m/s/gün)\\
\midrule
CoRoT-2\,b  & 1,743 & 12,5 & 0,151 & 26,7  & $1{,}7\times10^4$ & $-0{,}118$\\
HAT-P-37\,b & 2,797 & 13,4 & 0,108 & 1,7   & $6{,}8\times10^2$ & $-0{,}116$\\
KELT-16\,b  & 0,969 & 11,7 & 0,309 & 614,6 & $3{,}0\times10^5$ & $-0{,}273$\\
TrES-1\,b   & 3,030 & 11,4 & 0,101 & 0,8   & $7{,}2\times10^2$ & $-0{,}046$\\
TrES-5\,b   & 1,482 & 13,7 & 0,164 & 22,6  & $3{,}1\times10^4$ & $-0{,}062$\\
WASP-4\,b   & 1,338 & 12,5 & 0,188 & 29,6  & $6{,}2\times10^4$ & $-0{,}046$\\
WASP-10\,b  & 3,093 & 12,4 & 0,086 & 1,9   & $9{,}0\times10^2$ & $-0{,}085$\\
WASP-12\,b  & 1,091 & 11,6 & 0,336 & 452,5 & $2{,}0\times10^5$ & $-0{,}263$\\
WASP-22\,b  & 3,533 & 12,0 & 0,123 & 1,2   & $2{,}3\times10^2$ & $-0{,}193$\\
WASP-45\,b  & 3,126 & 11,8 & 0,100 & 1,1   & $8{,}7\times10^1$ & $-0{,}514$\\
XO-4\,b     & 4,125 & 10,8 & 0,122 & 3,3   & $7{,}1\times10^2$ & $-0{,}144$\\
HAT-P-7\,b  & 2,205 & 10,5 & 0,244 & 97,7  & (uzuyor) & $+0{,}079$\\
WASP-1\,b   & 2,520 & 11,3 & 0,176 & 11,0  & (uzuyor) & $+0{,}085$\\
WASP-46\,b  & 1,430 & 13,0 & 0,171 & 32,1  & (uzuyor) & $+0{,}343$\\
\bottomrule
\end{tabular}
\end{center}

\textbf{Sağlama:} WASP-12\,b için hesapladığım Q$'_\star = 2{,}0\times10^5$, Yee vd. (2020)'nin
$1{,}8\times10^5$ değeriyle uyumlu; WASP-4\,b için $6{,}2\times10^4$, Baştürk vd. (2025)'in $\sim8\times10^4$'üyle
uyumlu. Yani formül ve birimler doğru.

